\documentclass[10pt]{article}
\usepackage[letterpaper,margin=0.72in]{geometry}
\usepackage{fontspec}
\setmainfont[Path=/usr/share/texmf/fonts/opentype/public/lm/,UprightFont=lmroman10-regular.otf,BoldFont=lmroman10-bold.otf,ItalicFont=lmroman10-italic.otf,BoldItalicFont=lmroman10-bolditalic.otf]{lmroman10-regular.otf}
\usepackage{unicode-math}
\setmathfont[Path=/usr/share/texmf/fonts/opentype/public/lm-math/]{latinmodern-math.otf}
\usepackage{booktabs,tabularx,array,xcolor,enumitem,microtype,fancyhdr,titlesec,ragged2e}
\definecolor{darkgray}{gray}{0.28}
\pagestyle{fancy}
\fancyhf{}
\fancyhead[L]{Bethe (1928) -- mathematical notation audit}
\fancyhead[R]{Revision 4}
\fancyfoot[C]{\thepage}
\renewcommand{\headrulewidth}{0.3pt}
\setlength{\parindent}{0pt}
\setlength{\parskip}{0.42em}
\setlist[itemize]{leftmargin=1.35em,itemsep=0.22em,topsep=0.22em}
\titleformat{\section}{\large\bfseries}{}{0pt}{}
\newcommand{\pass}{\textbf{PASS -- source print confirmed}}
\begin{document}

\begin{center}
{\LARGE\bfseries Mathematical Notation and Formula Audit -- Revision 4}\\[0.35em]
{\large Hans Bethe, \textit{Theorie der Beugung von Elektronen an Kristallen} (1928)}\\[0.3em]
{\small Rechecked against the clearer 75-page \textit{Annalen der Physik} facsimile -- 17 September 2026}
\end{center}
\vspace{0.25em}\hrule\vspace{0.6em}

\textbf{Certification scope.} The mathematical layer of this English edition has now been checked page by page against the clearer supplied facsimile (original journal pp.~55--129). The audit covers all main equation numbers (1)--(94), all lettered and repeated numbered occurrences, displayed unnumbered formulas, inline mathematical notation, and the numerical/math tables. The rendered source page image, not the PDF text layer or OCR, is the authority for glyph decisions.

The companion audit ledger contains 158 numbered-equation occurrences and 2,661 indexed math-line/inline-math occurrences. A source-printed anomaly is therefore no longer an ``unresolved'' item: it receives a PASS when the English edition faithfully reproduces what Bethe actually printed, with an editorial note explaining the apparent source error.

\section{Five source-print anomalies rechecked at high resolution}

\begin{enumerate}[leftmargin=1.55em]
\item \textbf{Original p.~69, Eq.~(16) -- summation index.} \pass. The source clearly prints a summation subscript $g$ while the vacuum-wave amplitudes and propagation vectors in the summand carry the index $m$. On the next page the boundary-condition sums are over $m$. The printed $g$ is therefore very likely a typesetting error for $m$. The English body deliberately reproduces the printed $g$ and flags it rather than silently emending Bethe.

\item \textbf{Original pp.~80--81, Eq.~(41) -- missing square.} \pass. The source clearly prints
\[
|\gamma_h|<v_0.
\]
However Eq.~(41b) contains $\gamma_h^2$, and with $v_0=8\pi^2meV_0/h^2$ the derivation makes the dimensionally consistent condition $\gamma_h^2<v_0$. Thus Eq.~(41) almost certainly has a source typesetting error. The English equation remains faithful to the print; the audit note records the likely intended form.

\item \textbf{Original p.~94, Fig.~7 -- curve-label subscripts.} \pass. The prose immediately above Fig.~7, consistent with Eq.~(72), uses
\[
(\sigma_{11}\tau_{22}-\sigma_{12}\tau_{12})^2,
\]
whereas the label written on the solid curve in Fig.~7 is visibly
\[
(\sigma_{11}\tau_{22}-\sigma_{22}\tau_{11})^2.
\]
This is a real mismatch in the 1928 printing. The previous checklist described it imprecisely as a difference in ``reflection indices''; that wording is corrected here. The English edition reproduces each source location as printed.

\item \textbf{Original p.~102, Eq.~(37b) -- two internal formula anomalies.} \pass. The third line visibly has the denominator $2\varkappa^2-k_g^2$, while the final $V_{22}^{ii}$ numerator is visibly an unsquared $v_{g-h^i}$. Both differ from the pattern generated by Eq.~(35a): the weak-beam denominator there is $\varkappa^2-k_g^2$, and the $i=j$ specialization of the preceding product would give a square in the numerator. Both items are therefore very likely source typesetting errors. They are retained exactly as printed in the English transcription and explicitly documented.

\item \textbf{Original pp.~109, 117--118, Eq.~(92) -- Laguerre subscript.} \pass. Original p.~109 and Eq.~(91a) on p.~117 use $L_{n+l}^{2l+1}$, and the Waller integral quoted on p.~118 again uses $L_{n+l}^{2l+1}$. Eq.~(92) alone visibly prints $L_{n-l}^{2l+1}$. The isolated $n-l$ is therefore almost certainly a source typo. The English body reproduces Eq.~(92) as printed and includes a translator's note.
\end{enumerate}

\section{Corrections confirmed in the current body}

The following previously identified transcription errors have been corrected and were rechecked against the clearer scan:
\begin{itemize}
\item Fraktur propagation vector $\mathfrak f$ restored where earlier copies had $\mathfrak k$; Eqs.~(18)--(19) are confirmed as $\mathfrak f$ notation.
\item Original p.~92 before Eq.~(68): $2\pi g_n$ confirmed.
\item Original p.~92, Eq.~(70): the first $\tau_{22}$ denominator is $-4\pi(\mathfrak g,\mathfrak C_0-\pi\mathfrak g)$.
\item Original p.~103 after Eq.~(79): the German is \textit{anzunehmen}; ``assumed'' is correct. The earlier audit change to ``excluded'' was erroneous and has been reversed.
\item Original p.~104, Eq.~(80): $\pi$ restored in $\mathfrak C_0+\pi\mathfrak g$; Eq.~(81)'s final subscript is clearly $g_3$; Eq.~(82) uses Bethe's $\vartheta$ glyph.
\item Original pp.~109--117: the atomic/effective nuclear-charge symbol is italic $z$; the affected $x/\varkappa$-like misreadings have been corrected.
\item Original p.~121: the unnumbered asymptotic Fourier coefficient has denominator $(1+\alpha^2)^{2n}$.
\end{itemize}

\section{Status of this edition}

For the \textbf{mathematics and notation}, the current body is now certified against the supplied clear facsimile: every numbered formula and every indexed mathematical occurrence has been checked. The five items above are not remaining transcription failures; they are verified peculiarities of the 1928 source printing and are intentionally preserved with notes.

This mathematical certification does \emph{not} by itself claim that every prose sentence has completed the separate sentence-by-sentence bilingual translation audit.

\begin{center}
\fbox{\parbox{0.88\textwidth}{\centering
\textbf{Mathematical transcription errors found in the English copy have been corrected.}\\
\textbf{Verified source-print anomalies are preserved and identified, not silently modernized.}
}}
\end{center}

\vfill
{\small\color{darkgray}Editorial note: page references in this appendix are original journal page numbers.}
\end{document}
